% Nouveau EGA V (cote 25): archive cote 25, site page 163 (PDF page 164).
% typescript. Draft compared: none.
% AI-transcribed by Claude Opus 5.5 (claude-opus-5-5, adaptive thinking, effort high), Message Batches API, batch ega25b (msgbatch_01A7n5YuXubUPNK1zsepi8RG), 2026-11-13; not yet checked by eye.
\scanpage{25/163}
\mspage{\hand{6}}

un sous-préschéma X" de X, régulier, \var{\xout{\illegible{1}}}{}{typewriter strike-out}{ega25_25_163-v0001} contenant X, et tel que X
puisse se décrire dans X"\var{\xout{\illegible{1}}}{}{typewriter strike-out}{ega25_25_163-v0002} par une seule équation $\phi = 0$ , admettant x comme
zéro \var{\inserted{sing}}{sing}{insertion}{ega25_25_163-v0003} élémentaire de multiplicité n.

On conclut en particulier que si A est local noeth élémentaire de
\var{\xout{15.7 On peut donner des variantes géométriques de 15.\illegible{1} lorsque \illegible{1}}}{}{typewriter strike-out}{ega25_25_163-v0004}
mult. n, alors A est Coh. Mac. \var{\xout{\unsure{Il est clair}}}{}{typewriter strike-out}{ega25_25_163-v0005} Nous montrerons au par.20
\var{\xout{\illegible{2} (point de \unsure{x}, pas nécessairement rationnel \unsure{sur} k)}}{}{typewriter strike-out}{ega25_25_163-v0006}
(par un calcul facile d'éclatement) que le point fermé de Spec(A) est
le seul point singulier de Spec(A); il en résulte que \var{\xout{si dim}}{}{typewriter strike-out}{ega25_25_163-v0007} A est normal
sss dim A $\geqslant$ 2, réduit sss dimA $\geqslant$1. Si dim A = 0, alors  A est élémentaire
de multiplicité n sss  $m/m^2$ est de \var{\over{dim}{\unsure{rang}}}{\unsure{rang}}{writing over}{ega25_25_163-v0008}1 et  n est le plus petit entier
tel que $\underline{m}^n = 0$; on notera que ces anneaux sont aussi lesquotients d'anneaux
de valu\var{\xout{\unsure{at}}}{}{typewriter strike-out}{ega25_25_163-v0009}ion discrète par la\var{\inserted[typed]{\var{\xout{\illegible{1}}}{}{typewriter strike-out}{ega25_25_163-v0011}}}{\var{\xout{\illegible{1}}}{}{typewriter strike-out}{ega25_25_163-v0011}}{insertion}{ega25_25_163-v0010} puissance n.ème de \var{\over{l'}{leur}}{leur}{writing over}{ega25_25_163-v0012} idéal maximal.

15.7.   Il y a lieu de dégager les variantes ``géométriques'' de 15.6. ,
on trouvera en particulier que si X est un préschéma loc de type fini sur
le corps k, $x \in X$, et si \var{\over{n}{d}}{d}{writing over}{ega25_25_163-v0013} $= \dim_x X$ , alors  x est une singularité \var{\xout{\unsure{ordinaire}}}{}{typewriter strike-out}{ega25_25_163-v0014}
\var{\inserted[typed]{élémentaire}}{élémentaire}{insertion}{ega25_25_163-v0015}
\var{\xout{\illegible{1}}}{}{typewriter strike-out}{ega25_25_163-v0016} de multiplicité  n de X, sss il existe un voisinage ouvert U de x
tel que U se plonge comme sous-préschéma d'un préschéma lisse X' sur k,
\var{\xout{\illegible{1}}}{}{typewriter strike-out}{ega25_25_163-v0017} conne xe et de dimension \var{\over{n}{d}}{d}{writing over}{ega25_25_163-v0018}+1 , défini par une équation $\phi$=0 , x étant
un zéro \var{\inserted{géom.}}{géom.}{insertion}{ega25_25_163-v0019} singulier ordinaire de multiplicité n de $\phi$ . On aura dit d'ailleurs
\var{\margin{$\surd$}}{\placedmargin{$\surd$}}{marginal note placed}{ega25_25_163-v0020} dans 15.5. que cela signifie  que  la \var{\inserted{valeur en x de}}{valeur en x de}{insertion}{ega25_25_163-v0021} partie principale $d^n\phi$ de $\phi$ , \var{\xout{restrei}}{}{typewriter strike-out}{ega25_25_163-v0022}
\var{\xout{nte à X=V($\phi$)}}{}{typewriter strike-out}{ega25_25_163-v0023} (qui à priori est un\unsure{e} \var{\xout{\illegible{1}}}{}{typewriter strike-out}{ega25_25_163-v0024} \var{\inserted{élément}}{élément}{insertion}{ega25_25_163-v0025} de $\underline{P}^n_{X'/k}\otimes \unsure{\underline{k}(x)}$ ,et plus
précisément de son Idéal d'augmentation $\underline{I}$ )  est en fait un\var{\deleted{\unsure{e section}}}{}{deletion}{ega25_25_163-v0026} \var{\inserted{élément de}}{élément de}{insertion}{ega25_25_163-v0027}
\var{\xout{de}}{}{typewriter strike-out}{ega25_25_163-v0028} $\underline{I}^n = \underline{Sym}^n_{\underline{O}_U}(\underline{\Omega}^1_{X'/k}\otimes k(\hand{x}))$), et  à ce titre est une \var{\xout{\illegible{1}}}{}{typewriter strike-out}{ega25_25_163-v0029} \var{\inserted{forme}}{forme}{insertion}{ega25_25_163-v0030} lisse \var{\xout{en x}}{}{typewriter strike-out}{ega25_25_163-v0031}
\var{\margin{|}}{\placedmargin{|}}{marginal note placed}{ega25_25_163-v0032} \hand{[}NB Oubli dans 15.1. : l'ensemble des points de lissité d'une forme est
\var{\margin{|}}{\placedmargin{|}}{marginal note placed}{ega25_25_163-v0033} ouvert\hand{]}. On remarquera de même qu'un tel point est isolé dans l'ensemble
des points non lisses, il est géom normal \var{\xout{\illegible{1}}}{}{typewriter strike-out}{ega25_25_163-v0034} (resp. géom réduit) sss
\var{\over{n}{d}}{d}{writing over}{ega25_25_163-v0035} $\geqslant 2$  (resp. \var{\over{n}{d}}{d}{writing over}{ega25_25_163-v0036} $\geqslant 1$). Si \var{\over{n}{d}}{d}{writing over}{ega25_25_163-v0037} = 0 et k(x) = k (\var{\inserted{i.e. x  isolé et}}{i.e. x  isolé et}{insertion}{ega25_25_163-v0038} x rationnel sur k) alors \var{\deleted{x}}{}{deletion}{ega25_25_163-v0039}
point singulier élémentaire  de mult\hand{,} \unsure{n} signifie que $\underline{O}X,x$  est  k-\var{\deleted{\unsure{isomor}}}{}{deletion}{ega25_25_163-v0040}

\mspage{\hand{163}}
